Recursive AI-assisted software engineering systems increasingly operate through layered cycles of autonomous generation, evaluation, and refinement. As these recursive workflows scale, the cognitive burden required for humans to maintain contextual coherence grows alongside the probability of recursive drift — a progressive divergence between intended system behavior and the evolving semantic state of the system itself.

Existing software engineering methodologies were primarily designed around human-paced iteration, bounded cognitive context, and explicit synchronization boundaries. Recursive AI-assisted systems destabilize these assumptions by accelerating execution velocity, increasing recursive complexity, and continuously mutating contextual state across autonomous development cycles.

This paper introduces \textit{reflector}, a reflective development framework designed to reduce cognitive load while preserving human intentionality, governance, and architectural coherence within recursive AI-assisted workflows. \textit{reflector} introduces synchronization boundaries, reflective auditing systems, milestone-constrained recursive execution, and scoped autonomous agents as mechanisms for maintaining contextual stability across recursive development systems.

Rather than treating governance as an external control layer, \textit{reflector} frames synchronization as a contextual recompression process through which humans and AI systems periodically re-align around stabilized semantic state representations. We argue that reflective synchronization architectures will become increasingly necessary as recursive AI-assisted development systems continue scaling in autonomy, complexity, and execution speed.